\chapter{Diseases of the Intestines}
\label{ch:intestines}
The intestines are divided into the small intestine (duodenum, jejunum, and ileum) and the large intestine (cecum, colon, rectum, and anal canal). Together, they span approximately 7.5 meters and are responsible for nutrient absorption, water reclamation, and waste elimination. Intestinal diseases include inflammatory conditions, infections, neoplasms, motility disorders, and vascular emergencies.

%-------------------------------------------------------------
\section{Acute Abdominal Emergencies}
%-------------------------------------------------------------

%-------------------------------------------------------------

%-------------------------------------------------------------

%-------------------------------------------------------------

%-------------------------------------------------------------

%-------------------------------------------------------------

%-------------------------------------------------------------

%-------------------------------------------------------------

\label{sec:Acute Abdominal Emergencies}
%-------------------------------------------------------------%-------------------------------------------------------------

\subsection{Appendicitis}
\index{Appendicitis}
\label{sec:Appendicitis}
Appendicitis is inflammation of the vermiform appendix, the most common surgical emergency of the abdomen, with a lifetime risk of approximately 7--8\%. It typically occurs between ages 10 and 30, with a slight male predominance. The condition usually results from luminal obstruction of the appendix by a fecalith, lymphoid hyperplasia, or rarely tumors, leading to increased intraluminal pressure, bacterial overgrowth, ischemia, and perforation. Clinical features include the classic presentation of periumbilical pain that migrates to the right lower quadrant (McBurney's point), accompanied by anorexia, nausea, vomiting, and low-grade fever; however, the classic presentation is present in only 50--60\% of cases. Physical findings include tenderness at McBurney's point, guarding, and rebound tenderness. The Alvarado score (MANTRELS) aids clinical assessment. Diagnosis involves laboratory findings of leukocytosis with left shift and CT abdomen/pelvis as the preferred imaging modality in adults (sensitivity and specificity >95\%); ultrasound may be used as first-line imaging in children and pregnant women. Management is appendectomy (laparoscopic preferred over open), with preoperative IV antibiotics; antibiotics alone may be considered for uncomplicated appendicitis in selected patients, though recurrence rates of 15--40\% exist. Complicated appendicitis (gangrenous, perforated, with abscess or peritonitis) requires IV antibiotics and possibly interval appendectomy. Prognosis is excellent with timely surgical intervention.

%-------------------------------------------------------------
\section{Celiac Disease}
%-------------------------------------------------------------

%-------------------------------------------------------------

%-------------------------------------------------------------

%-------------------------------------------------------------

%-------------------------------------------------------------

%-------------------------------------------------------------

%-------------------------------------------------------------

%-------------------------------------------------------------

\label{sec:section:Celiac Disease}
%-------------------------------------------------------------%-------------------------------------------------------------

\subsection{Celiac Disease}
\index{Celiac Disease}
\label{sec:Celiac Disease}
Celiac disease is an autoimmune enteropathy triggered by ingestion of gluten (gliadin and glutenin proteins found in wheat, barley, and rye) in genetically susceptible individuals (HLA-DQ2 or HLA-DQ8). Prevalence is approximately 1\% worldwide. The condition results from an immune response causing villous atrophy, crypt hyperplasia, and intraepithelial lymphocytosis in the small intestinal mucosa, primarily affecting the duodenum and jejunum. Clinical features vary widely: classical symptoms in children include chronic diarrhea, failure to thrive, and abdominal distension; in adults, diarrhea, weight loss, iron deficiency anemia, and osteoporosis. Dermatitis herpetiformis is the cutaneous manifestation. Many patients are asymptomatic or have extraintestinal manifestations (fatigue, neurological symptoms, liver disease, infertility). Diagnosis requires serological testing (anti-tissue transglutaminase IgA, anti-endomysial antibodies) and duodenal biopsy showing villous atrophy, crypt hyperplasia, and increased intraepithelial lymphocytes; HLA testing helps exclude the diagnosis. Management is a strict lifelong gluten-free diet, which leads to mucosal recovery and symptom resolution. Prognosis is excellent with dietary compliance.

%-------------------------------------------------------------
\section{Colorectal Neoplasms}
%-------------------------------------------------------------

%-------------------------------------------------------------

%-------------------------------------------------------------

%-------------------------------------------------------------

%-------------------------------------------------------------

%-------------------------------------------------------------

%-------------------------------------------------------------

%-------------------------------------------------------------

\label{sec:Colorectal Neoplasms}
%-------------------------------------------------------------%-------------------------------------------------------------

\subsection{Colorectal Adenoma}
\index{Colorectal Adenoma}
\label{sec:Colorectal Adenoma}
Colorectal adenoma is a benign epithelial neoplasm that is a precursor to colorectal cancer through the adenoma-carcinoma sequence. It is classified as tubular, tubulovillous, or villous based on histology, and by the degree of dysplasia (low-grade or high-grade). Risk factors include age (>50), male sex, family history, hereditary polyposis syndromes (FAP, Lynch syndrome), inflammatory bowel disease, obesity, smoking, and high red meat consumption. Most adenomas are asymptomatic and detected through screening colonoscopy. Larger adenomas (>1 cm), villous histology, and high-grade dysplasia carry higher malignant potential. Diagnosis is based on colonoscopy with histopathological evaluation. Management involves endoscopic removal (polypectomy) with surveillance colonoscopy intervals based on polyp number, size, and histology. Prognosis is excellent with complete removal.

\subsection{Colorectal Cancer}
\index{Colorectal Cancer}
\label{sec:Colorectal Cancer}
Colorectal cancer (CRC) is the third most commonly diagnosed cancer worldwide and the second leading cause of cancer-related death. Approximately 70\% of CRCs arise from adenomatous polyps via the chromosomal instability pathway (APC/KRAS/TP53 mutations); Lynch syndrome (hereditary nonpolyposis colorectal cancer, HNPCC) accounts for 2--4\% of cases due to mismatch repair gene mutations. The condition results from accumulated genetic and epigenetic alterations driving malignant transformation. Clinical features vary by location; right-sided colon cancers tend to present with iron deficiency anemia and occult GI bleeding, while left-sided cancers present with altered bowel habits, obstruction, and hematochezia. Diagnosis is based on colonoscopy with biopsy; screening methods include colonoscopy (gold standard), fecal occult blood testing (FIT), and CT colonography. Staging uses TNM classification with CT, MRI (rectal cancer), and PET. Management of localized disease is surgical resection with lymphadenectomy, with adjuvant chemotherapy (FOLFOX, CapeOX) for stage III disease; neoadjuvant chemoradiation is standard for locally advanced rectal cancer. Metastatic CRC is treated with systemic chemotherapy, targeted therapy (bevacizumab [anti-VEGF], cetuximab/panitumumab [anti-EGFR] for RAS wild-type tumors), immunotherapy (pembrolizumab for MSI-high/dMMR tumors), and metastasectomy for selected patients. Prognosis depends on stage at diagnosis.

\subsection{Familial Adenomatous Polyposis}
\index{Familial Adenomatous Polyposis}
\label{sec:Familial Adenomatous Polyposis}
Familial adenomatous polyposis (FAP) is an autosomal dominant condition caused by germline mutations in the APC gene on chromosome 5q21, characterized by the development of hundreds to thousands of adenomatous polyps throughout the colon and rectum, typically beginning in adolescence. The condition results from APC gene mutation leading to uncontrolled cell proliferation. Clinical features include the development of polyps in adolescence; without treatment, CRC develops in virtually 100\% of patients by age 40. Diagnosis is based on colonoscopy showing innumerable polyps and genetic testing. Management involves prophylactic proctocolectomy with ileal pouch-anal anastomosis (IPAA) recommended in the late teens to early twenties. Prognosis is excellent with prophylactic surgery; attenuated FAP (AFAP) has fewer polyps (10--100) and later-onset CRC. Gardner syndrome is a variant with extraintestinal manifestations (osteomas, desmoid tumors, congenital hypertrophy of retinal pigment epithelium).

\subsection{Lynch Syndrome}
\index{Lynch Syndrome}
\label{sec:Lynch Syndrome}
Lynch syndrome (hereditary nonpolyposis colorectal cancer, HNPCC) is an autosomal dominant cancer predisposition syndrome caused by germline mutations in DNA mismatch repair genes (MLH1, MSH2, MSH6, PMS2) or EPCAM deletions. It accounts for 2--4\% of all CRCs, with a lifetime CRC risk of 40--80\%. Associated cancers include endometrial (highest risk after CRC), ovarian, gastric, urinary tract, hepatobiliary, and small bowel cancers. The condition results from defective mismatch repair leading to microsatellite instability (MSI) or mismatch repair deficiency (dMMR). Clinical features include early-onset colorectal cancer and increased risk of multiple primary cancers. Diagnosis is based on Amsterdam criteria, revised Bethesda guidelines, immunohistochemistry of tumor tissue, and germline genetic testing. Management includes surveillance colonoscopy starting at age 20--25 (or 2--5 years before the earliest CRC in the family) every 1--2 years; prophylactic surgery may be considered. Prognosis is improved with surveillance; immunotherapy with pembrolizumab is particularly effective for MSI-high/dMMR cancers.

%-------------------------------------------------------------
\section{Diverticular Disease}
%-------------------------------------------------------------

%-------------------------------------------------------------

%-------------------------------------------------------------

%-------------------------------------------------------------

%-------------------------------------------------------------

%-------------------------------------------------------------

%-------------------------------------------------------------

%-------------------------------------------------------------

\label{sec:Diverticular Disease}
%-------------------------------------------------------------%-------------------------------------------------------------

\subsection{Diverticulitis}
\index{Diverticulitis}
\label{sec:Diverticulitis}
Diverticulitis is inflammation or infection of a diverticulum, occurring when a diverticulum becomes obstructed by fecal material, leading to microperforation and localized inflammation. Clinical features of uncomplicated diverticulitis include left lower quadrant pain, fever, leukocytosis, and altered bowel habits; complicated diverticulitis may involve abscess formation, perforation with peritonitis, fistula formation (colovesical fistula most common), or stricture. Diagnosis is based on CT abdomen/pelvis with contrast, which demonstrates pericolonic fat stranding, wall thickening, and abscess. Management of uncomplicated diverticulitis includes antibiotics (fluoroquinolones + metronidazole, or amoxicillin-clavulanate) and a liquid diet progressing to a high-fiber diet; complicated diverticulitis requires hospitalization, IV antibiotics, and percutaneous drainage for abscesses. Prognosis is good with appropriate treatment; elective colectomy is considered after two or more episodes of complicated diverticulitis or for immunocompromised patients.

\subsection{Diverticulosis}
\index{Diverticulosis}
\label{sec:Diverticulosis}
Diverticulosis is the presence of small, bulging pouches (diverticula) in the wall of the colon, most commonly in the sigmoid colon. It affects over 50\% of individuals over 60 in Western countries. The condition results from increased intraluminal pressure at weak points in the colonic wall where vasa recta penetrate the circular muscle layer, and is strongly associated with low-fiber diets, obesity, and aging. Clinical features include mostly asymptomatic presentation, discovered incidentally on colonoscopy or CT. Diagnosis is based on colonoscopy or imaging. Management includes high-fiber diets and stool softeners for prevention of complications. Prognosis is excellent; most patients remain asymptomatic.

%-------------------------------------------------------------
\section{Functional Gastrointestinal Disorders}
%-------------------------------------------------------------

%-------------------------------------------------------------

%-------------------------------------------------------------

%-------------------------------------------------------------

%-------------------------------------------------------------

%-------------------------------------------------------------

%-------------------------------------------------------------

%-------------------------------------------------------------

\label{sec:Functional Gastrointestinal Disorders}
%-------------------------------------------------------------%-------------------------------------------------------------

\subsection{Irritable Bowel Syndrome}
\index{Irritable Bowel Syndrome}
\label{sec:Irritable Bowel Syndrome}
Irritable bowel syndrome (IBS) is a chronic functional gastrointestinal disorder characterized by recurrent abdominal pain associated with altered bowel habits, in the absence of structural or biochemical abnormalities. It affects approximately 10--15\% of the global population and is more common in women. The condition results from visceral hypersensitivity, altered gastrointestinal motility, gut-brain axis dysfunction, intestinal dysbiosis, low-grade mucosal inflammation, and post-infectious triggers. IBS is subtyped by predominant stool pattern: IBS with constipation (IBS-C), IBS with diarrhea (IBS-D), IBS with mixed bowel habits (IBS-M), and unsubtyped IBS (IBS-U). Clinical features include recurrent abdominal pain, on average at least 1 day per week in the last 3 months, associated with defecation, change in stool frequency, or change in stool form. Diagnosis is based on the Rome IV criteria; alarm features (rectal bleeding, nocturnal symptoms, unintentional weight loss, family history of colorectal cancer, iron deficiency anemia, onset after age 50) warrant further investigation to exclude organic disease. Management is multimodal: dietary modification (low FODMAP diet, soluble fiber supplementation), antispasmodics (hyoscine, dicyclomine, peppermint oil), laxatives for IBS-C (linaclotide, lubiprostone, plecanatide), antidiarrheals for IBS-D (loperamide, eluxadoline, rifaximin), neuromodulators (tricyclic antidepressants, SSRIs), and gut-directed psychological therapies (hypnotherapy, cognitive-behavioral therapy). Prognosis is variable; symptoms are chronic but can be managed with multimodal therapy.

%-------------------------------------------------------------
\section{Inflammatory Bowel Disease}
%-------------------------------------------------------------

%-------------------------------------------------------------

%-------------------------------------------------------------

%-------------------------------------------------------------

%-------------------------------------------------------------

%-------------------------------------------------------------

%-------------------------------------------------------------

%-------------------------------------------------------------

\label{sec:Inflammatory Bowel Disease}
%-------------------------------------------------------------%-------------------------------------------------------------

\subsection{Crohn's Disease}
\index{Crohn's Disease}
\label{sec:Crohn's Disease}
Crohn's disease (CD) is a chronic, relapsing inflammatory bowel disease that can affect any part of the gastrointestinal tract from mouth to anus, though the terminal ileum and colon are most commonly involved. Peak incidence occurs in the second to fourth decades of life, with a bimodal distribution. The condition results from an interplay of genetic susceptibility (over 240 identified loci, with NOD2/CARD15 being the most significant), environmental factors (smoking doubles the risk), and dysregulated immune response, characterized by transmural inflammation, skip lesions, and non-caseating granulomas. Clinical features include chronic diarrhea (often non-bloody), abdominal pain (typically right lower quadrant), weight loss, fatigue, and perianal disease (fistulas, skin tags, abscesses); extraintestinal manifestations include arthritis, erythema nodosum, pyoderma gangrenosum, uveitis, and primary sclerosing cholangitis. Diagnosis involves colonoscopy with biopsies, MR enterography or CT enterography, and serological markers (CRP, fecal calprotectin). Management follows a step-up or top-down approach: 5-aminosalicylates for mild colonic disease, corticosteroids for moderate flares (not for maintenance), immunomodulators (azathioprine, methotrexate), and biologic therapies (anti-TNF [infliximab, adalimumab], anti-integrin [vedolizumab], anti-IL12/23 [ustekinumab], anti-IL23 [risankizumab]). Prognosis is variable; surgery is required in approximately 50\% of patients within 10 years but is not curative.

\subsection{Indeterminate Colitis}
\index{Indeterminate Colitis}
\label{sec:Indeterminate Colitis}
Indeterminate colitis is inflammatory bowel disease confined to the colon where features of both Crohn's disease and ulcerative colitis are present, and definitive classification is not possible even after pathological review. It accounts for 5--15\% of IBD cases. The condition results from overlapping pathological features of both Crohn's disease and ulcerative colitis. Clinical features are those of colonic inflammation without clear distinguishing characteristics. Diagnosis is based on histopathology that cannot definitively classify the disease as either CD or UC. Management follows the same principles as UC and CD, with treatment tailored to the predominant clinical phenotype. Prognosis is variable and depends on the clinical course.

\subsection{Ulcerative Colitis}
\index{Ulcerative Colitis}
\label{sec:Ulcerative Colitis}
Ulcerative colitis (UC) is a chronic inflammatory bowel disease limited to the colon and rectum, characterized by continuous mucosal and submucosal inflammation beginning at the rectum and extending proximally in a continuous fashion. The condition results from a dysregulated immune response to gut microbiota in genetically susceptible individuals; unlike Crohn's disease, it does not involve the small intestine (except for backwash ileitis in pancolitis) and does not cause transmural disease, strictures, or fistulas. Clinical features include bloody diarrhea, urgency, tenesmus, and abdominal cramping; extraintestinal manifestations similar to Crohn's disease include arthritis, erythema nodosum, and primary sclerosing cholangitis (occurring in 2--5\% of UC patients). Diagnosis is based on colonoscopy with biopsies showing continuous inflammation, crypt abscesses, and mucosal ulceration; severity is classified as mild, moderate, or severe based on stool frequency, rectal bleeding, hemoglobin, and CRP. Management includes 5-aminosalicylates (mesalamine) for mild-moderate disease, corticosteroids for flares, immunomodulators, and biologics (anti-TNF, vedolizumab, ustekinumab, tofacitinib [JAK inhibitor]); acute severe UC is a medical emergency treated with IV corticosteroids, with colectomy for refractory disease. Prognosis is variable; total proctocolectomy is curative.

%-------------------------------------------------------------
\section{Intestinal Obstruction}
%-------------------------------------------------------------

%-------------------------------------------------------------

%-------------------------------------------------------------

%-------------------------------------------------------------

%-------------------------------------------------------------

%-------------------------------------------------------------

%-------------------------------------------------------------

%-------------------------------------------------------------

\label{sec:Intestinal Obstruction}
%-------------------------------------------------------------%-------------------------------------------------------------

\subsection{Large Bowel Obstruction}
\index{Large Bowel Obstruction}
\label{sec:Large Bowel Obstruction}
Large bowel obstruction (LBO) is a blockage of the large intestine, less common than SBO. The most common cause in Western countries is colorectal cancer (50--60\%), followed by diverticulitis strictures and volvulus (sigmoid volvulus in elderly/bedridden patients, cecal volvulus in younger patients). The condition results from mechanical or functional obstruction of the colon. Clinical features include progressive abdominal distension, constipation, and eventually nausea and vomiting (late in the course). Diagnosis is based on abdominal X-ray showing a dilated colon with a cutoff point; CT abdomen/pelvis is diagnostic. Management depends on the cause; sigmoid volvulus may be decompressed by sigmoidoscopy; malignant LBO typically requires surgical decompression (colostomy or resection) as a palliative or curative measure depending on disease stage. Prognosis depends on the underlying cause and presence of complications.

\subsection{Small Bowel Obstruction}
\index{Small Bowel Obstruction}
\label{sec:Small Bowel Obstruction}
Small bowel obstruction (SBO) is a partial or complete blockage of the small intestine. The most common cause is adhesions from prior abdominal surgery (60\% of cases); other causes include hernias (inguinal, incisional), malignancy, volvulus, intussusception (in children), inflammatory bowel disease strictures, and gallstone ileus. The condition results from mechanical blockage of intestinal contents. Clinical features include crampy abdominal pain (periumbilical), nausea, vomiting (bilious early, feculent late), abdominal distension, and obstipation. Diagnosis is based on abdominal X-ray showing air-fluid levels and dilated proximal bowel loops with collapsed distal bowel; CT abdomen/pelvis is more sensitive and identifies the transition point and cause. Management includes nasogastric decompression, IV fluid resuscitation, electrolyte correction, and close monitoring; complete obstruction or signs of strangulation (fever, tachycardia, peritoneal signs, metabolic acidosis) require emergent surgical exploration. Prognosis depends on cause and presence of strangulation.

%-------------------------------------------------------------
\section{Short Bowel Syndrome}
%-------------------------------------------------------------

%-------------------------------------------------------------

%-------------------------------------------------------------

%-------------------------------------------------------------

%-------------------------------------------------------------

%-------------------------------------------------------------

%-------------------------------------------------------------

%-------------------------------------------------------------

\label{sec:section:Short Bowel Syndrome}
%-------------------------------------------------------------%-------------------------------------------------------------

\subsection{Anal Fissure}
\index{Anal Fissure}
\label{sec:Anal Fissure}
Anal fissure is a painful linear tear or ulceration in the distal anal canal lining extending from the dentate line to the anal verge. The condition results from trauma caused by passage of hard, dry stool or severe diarrhea, precipitating internal anal sphincter spasm and ischemia of the posterior anal midline (where 90\% of fissures occur due to poor microvascular perfusion). Clinical features include severe, sharp, tearing pain during and after defecation, accompanied by bright red rectal bleeding on toilet paper or stool surface. Physical examination reveals a visible longitudinal fissure; chronic fissures feature a sentinel skin tag and hypertrophied anal papilla. Diagnosis is clinical. Management includes stool softeners, high-fiber diet, warm sitz baths, and topical nitrates (0.2\% nitroglycerin) or calcium channel blockers (2\% diltiazem) to relax the sphincter; lateral internal sphincterotomy is reserved for chronic medically refractory fissures.

\subsection{Short Bowel Syndrome}
\index{Short Bowel Syndrome}
\label{sec:Short Bowel Syndrome}
Short bowel syndrome (SBS) is a malabsorptive condition resulting from surgical resection of a significant portion of the small intestine, leading to insufficient nutrient and fluid absorption. The most common cause is Crohn's disease, followed by mesenteric ischemia, trauma, and malignancy. The condition results from reduced intestinal absorptive surface area; the remaining functional bowel length and the presence or absence of the ileum and colon determine severity. Clinical features include profuse diarrhea, steatorrhea, weight loss, dehydration, electrolyte imbalances, and nutritional deficiencies. Diagnosis is based on surgical history and clinical presentation; adaptation of the remaining bowel occurs over months to years. Management includes parenteral nutrition and fluids, gradual enteral feeding to promote adaptation, antidiarrheal agents, and H2 blockers or PPIs to reduce gastric acid; teduglutide (a GLP-2 analog) promotes intestinal adaptation and may reduce parenteral nutrition dependence. Prognosis depends on the extent of resection and degree of adaptation; intestinal transplantation is reserved for patients with life-threatening complications of parenteral nutrition.

